% !TEX root = ../main.tex
\chapter{练习题答案线索与评分要点}
\label{ch:exercise-rubrics}

本章给出各章练习的最低得分点。基础题每题 4 分，中等题每题 6 分，困难题每题 8 分；答到下列关键点可得约 70\% 分数，余分取决于推导完整性、量纲、前提和反例。开放题允许不同方案，但必须明确模型层级与可验证量。

\section{第 1 章}
\begin{enumerate}
  \item 大面积降低衍射角并提高可用功率，却增加横向模密度和热/注入不均匀；单模要求目标模对所有竞争模保持净增益裕量。
  \item 边发射激光器依靠纵向 Fabry--P\'erot/DFB 反馈；VCSEL 依靠上下 DBR 的垂直反馈；PCSEL 依靠二维周期散射形成面内相干反馈并耦合到法向通道。
  \item 合格反例：只有表面光栅出光、腔反馈仍完全由边缘镜面或垂直 DBR 提供的器件；理由是二维光子晶体没有承担主要相干反馈。
  \item 高 $Q$ 只约束被动总衰减；阈值电流还取决于有源区重叠、内部损耗、载流子输运、热回写和竞争模。
\end{enumerate}

\section{第 2 章}
\begin{enumerate}
  \item 代入 $\ee^{-\ii\tilde\omega t}$：$\tilde\omega=\omega_0-\ii\gamma$ 给出 $\ee^{-\ii\omega_0t}\ee^{-\gamma t}$，故场幅按 $\gamma$ 衰减、能量按 $2\gamma$ 衰减。
  \item Bloch 模属于无限周期本征问题；被动共振模还包含开放边界与复频率；激射模由增益、饱和、注入、热和模式竞争共同选出。
  \item 示例：$g_{\mathrm{mat,th}}$ 低，但接触电阻大、电流扩展差或非辐射复合强，使达到该载流子密度所需电流仍很高。
\end{enumerate}

\section{第 3 章}
\begin{enumerate}
  \item 对无源介质由 $\nabla\times\E=\ii\omega\mu\Hf$ 消去 $\E$，得 $\nabla\times\varepsilon^{-1}\nabla\times\Hf=(\omega/c)^2\mu\Hf$。磁场形式可避免某些散度伪模；电场形式更便于介电界面和电偶极源，但两者都需正确函数空间。
  \item 面发射意味着能量穿过有限计算域进入自由空间，故必须采用出射辐射条件、端口或 PML，而非封闭边界。
  \item 由 $U(t)=U_0\ee^{-2\gamma t}$、$P_{\mathrm{loss}}=2\gamma U$ 与 $Q=\omega_0U/P_{\mathrm{loss}}$，得 $Q=\omega_0/(2\gamma)$。
  \item 出射 QNM 在复频率下具有径向因子 $\ee^{\ii\tilde k r}$；被动模的 $\Im\tilde k<0$ 使其在空间无穷远增长，普通 $L^2$ 积分发散，需 QNM/PML 正规化。
  \item 示例：PML 过近产生人工反射，会把边界腔误判为高 $Q$ 模，并改变远场和模式排序；答案需给出至少一个随 PML 扫描可识别的症状。
\end{enumerate}

\section{第 4 章}
\begin{enumerate}
  \item 方格可取 $\bm a_1=a\hat x,\bm a_2=a\hat y$，$\bm b_1=2\pi\hat x/a,\bm b_2=2\pi\hat y/a$；三角格可取 $a(1,0)$、$a(1/2,\sqrt3/2)$ 并由 $\bm a_i\cdot\bm b_j=2\pi\delta_{ij}$ 求倒格矢。
  \item 周期平移算符彼此对易且与 Maxwell 算符对易，公共本征态满足 $T_{\bm R}\psi=\lambda_{\bm R}\psi$；群乘法要求 $\lambda_{\bm R}=\ee^{\ii\kvec\cdot\bm R}$。
  \item 因 $\ee^{\ii(\kvec+\gvec)\cdot\bm R}=\ee^{\ii\kvec\cdot\bm R}$，其中 $\gvec\cdot\bm R=2\pi m$，故两者具有同一 Bloch 边界条件。
  \item $\Gamma$ 点对应法向相位匹配，并有最高 little-group 对称性和多支基本波简并；辐射矩阵元仍可能因对称性为零，不能只看群速度。
  \item 失效例：上下层栈强不对称、深刻蚀或高折射率反差导致 $E_z/H_z$ 与面内分量强混合，此时需三维全矢量模型。
\end{enumerate}

\section{第 5 章}
\begin{enumerate}
  \item 参数是 Bloch 波矢 $\kvec$，对每个 $\kvec$ 求 Maxwell 算符本征值 $\omega_n(\kvec)$；整组本征值随参数组成能带。
  \item 电子能带给出载流子允许能量、占据和跃迁；光子能带给出电磁模式允许频率、对称性与传播。
  \item 高对称路径只是一维采样；路径外可能存在模式填入视觉空白，真正完全带隙需检查整个不可约布里渊区。
  \item 带边是色散极值，$\Gamma$ 点是倒空间位置，激射模还需开放损耗、增益、注入、热和非线性竞争；三者定义层级不同。
  \item 至少检查辐射通道/复频率、量子阱与损耗层重叠、有限阵列尺寸收敛、电流分布和热回写中的三项。
\end{enumerate}

\section{第 6 章}
\begin{enumerate}
  \item 同一倒格矢集合既可把 $+\beta$ 散射到 $-\beta$ 形成反馈，也可把导模分量移到 $k_\parallel\simeq0$ 打开法向辐射；两者由不同 Fourier 分量/阶次控制。
  \item 辐射增强后应检查总阈值是否上升、外耦合效率是否改善，以及竞争模排序/远场/偏振是否被改写，三类任取三项并说明方向不必相同。
  \item 光锥外分量没有传播型自由空间通道；即使 slab 中存在，若周期散射未把它移入光锥或矩阵元为零，法向辐射仍弱。
  \item 合格场景需满足 $\Gamma$ 点附近由 $\pm x,\pm y$ 四个主导波构成、其他 Fourier 阶次远失谐、层栈可压缩为少数参数，并说明何时升级模型。
\end{enumerate}

\section{第 7 章}
\begin{enumerate}
  \item 先由 $Q=\omega_0/(2\gamma)$ 得时间能量损耗率 $2\gamma$；再用能量速度 $v_{\mathrm{en}}$ 得空间功率损耗 $\alpha=2\gamma/v_{\mathrm{en}}=\omega_0/(Qv_{\mathrm{en}})$。
  \item 模增益是材料增益按模式能量密度在有源区加权的结果；场分布不同导致 $\Gamma_{\mathrm{opt}}$、谱失谐和空间载流子权重不同。
  \item 示例：提高约束因子同时把场推入高吸收掺杂层，或需要更差的电流扩展/更高温升，净损耗增加可使 $I_{\mathrm{th}}$ 上升。
  \item 式中 $g_{\mathrm{mat}}[N(\rvec),T(\rvec)]$ 与模场权重相乘，正是电学/热学输出进入光学阈值的接口；均匀载流子近似不成立时不可用单一 $\Gamma$ 代替。
\end{enumerate}

\section{第 8 章}
\begin{enumerate}
  \item $\Gamma$ 点的波矢在 $C_{4v}$ 全部操作下保持不变，故 little group 为 $C_{4v}$；二维 irrep $E$ 的两个基函数在理想结构中简并，可对应一对正交偏振候选。
  \item 法向远场是单胞电流/场的相干积分；奇对称分量在中心通道两侧相消，故中心可出现零点，但需说明具体端口表示。
  \item $C_{6v}$ 法向矢量通道按 $(x,y)$ 变换，属于 $E_1$；$E_2$ 与其正交，理想 $\Gamma$ 点耦合可被对称性禁止。
  \item 场景示例：两偏振频率近简并但金属吸收重叠不同；或上下辐射/边缘泄漏对两偏振不同。答案需比较总损耗而非只比较频率。
  \item 上下不对称改变两个半空间的端口阻抗与相位，使上/下辐射振幅、偏振椭圆和角谱干涉都改变，不只是总功率。
  \item 一条可行链：电流升高 $\to$ 非均匀温升/载流子折射率变化 $\to$ 对称性降低和两偏振净增益差缩小或翻转 $\to$ 偏振混合/切换。
\end{enumerate}

\section{第 8B 章}
\begin{enumerate}
  \item BIC 的频率位于至少一个传播辐射通道的连续谱中，但对所有开放通道的耦合振幅为零；封闭腔没有这一“嵌入连续谱”定义。
  \item 模式与连续谱处于同一频段只说明通道存在；实际辐射还需非零耦合矩阵元，矩阵元可因对称或多通道干涉相消。
  \item 方格晶格法向电场通道按 $(x,y)$ 变换为 $E$；一维 irrep 与该端口表示正交时，理想结构可形成对称保护 BIC。需同时考虑上下半空间与所有偏振通道。
  \item 若 $Q\propto\delta^{-2}$，接近理想点会降低辐射损耗，却使相对几何误差、有限尺寸和无序决定实际泄漏；输出耦合也可能过弱。
  \item 开放题按“扰动降低何种对称性—新允许哪个端口矩阵元—首先改变何种观测量”评分；需给出明确孔形或位移。
\end{enumerate}

\section{第 9 章}
\begin{enumerate}
  \item 二阶 Bragg 条件下扩展区的 $\pm\bm G_x,\pm\bm G_y$ 导模分量折叠到 $\Gamma$，它们近简并且由最低阶相关 Fourier 分量互耦；任选四个平面波不满足这一物理来源。
  \item 四向量应各自归一且两两内积为零；在 $C_{4v}$ 操作下分别闭合为 $A_1/B_1/E$ 子空间。答案需至少展示一个内积和一个群操作。
  \item 孔半径主要改变平面 Fourier 系数；刻蚀深度改变调制强度与纵向重叠；层栈厚度改变 $\phi_0(z)$、失谐和重叠，三者均非单一比例缩放。
  \item 保持 $C_{4v}$ 时不同 irrep 的矩阵元由 Schur 正交性为零；降到较低子群后原 irrep 可分解到同一子群表示，从而允许混合。
  \item 该式只求小信号模增益补偿光学损耗；阈值电流还需把所需材料增益映射到载流子分布和电流，并含热回写。
  \item 失效例需指出缺失项，如高阶 Bloch 波、多个纵模、强开放通道、有限边缘或空间非均匀增益，并说明可用 3D CWT/RCWA/FDTD/FEM 中哪一种补上。
\end{enumerate}

\section{第 10 章}
\begin{enumerate}
  \item 示例：多层 slab 中两个纵模都接近 Bragg 共振，二维四波模型缺失纵向混合，而保留多纵模与高阶波的 3D CWT 仍可有效。
  \item $m$ 标记面内倒格通道，$n$ 标记纵向本征函数；多个 $n$ 近简并或具有不同辐射/增益重叠时，一个 $n_{\mathrm{eff}}$ 无法保留其耦合矩阵。
  \item 开放通道自能的实部给频移、反厄米部分给衰减；二者同源于同一 Green 函数并受因果性联系。
  \item 吸收环提高高阶包络在边缘的损耗，增大模鉴别；目标模也有非零边缘尾，故其总损耗和阈值同时上升。
\end{enumerate}

\section{第 11 章}
\begin{enumerate}
  \item CWT 从 Maxwell 方程与周期 Fourier 分量构造传播波耦合，保留空间结构；TCMT 以少数共振振幅和端口为状态，强调输入输出。前者解释几何耦合来源，后者解释开放端口响应。
  \item 非厄米算符的右模不按普通共轭内积正交；用左模投影才能消去其他模并得到正确分母，例外是具特殊对称/正规性的算符。
  \item 厄米弱扰动下大场区贡献大；开放系统还受左/右模相位、材料张量、辐射通道和制造约束影响，最大 $|E|$ 处不保证目标频移/损耗比最优。
  \item 当两模实频差小而衰减率/非厄米耦合差大时，先达到阈值由损耗分裂主导；答案应从复本征值虚部或净增益比较说明。
  \item 可行链：温升改变折射率或各向异性 $\to$ 两偏振的对角失谐/非厄米耦合变化 $\to$ 复本征值越过或排序翻转 $\to$ 偏振切换。
\end{enumerate}

\section{第 12 章}
\begin{enumerate}
  \item 将 $H_z=\sum_{\gvec}H_{\gvec}\ee^{\ii(\kvec+\gvec)\cdot\rvec}$ 与 $\eta=\sum\eta_{\gvec}\ee^{\ii\gvec\cdot\rvec}$ 代入标量方程；乘积产生卷积索引 $\eta_{\gvec-\gvec'}$，投影到每个平面波即得矩阵。
  \item 对圆盘积分 $A_c^{-1}\int_{r<R}\ee^{-\ii\gvec\cdot\rvec}\dd^2r$；角积分给 $2\pi J_0(Gr)$，再用 $\int rJ_0(Gr)\dd r=RJ_1(GR)/G$，得填充因子乘 $2J_1(GR)/(GR)$。
  \item 归一化频率 $\Omega=a/\lambda$，故 $a=\Omega_0\times\SI{940}{nm}$；不换算就无法把无量纲带图与材料增益波段、工艺尺寸对应。
  \item PWEM 平带不含开放辐射、内部损耗、有限边缘、增益重叠和热/注入；至少列三项并指出需哪个后续模型。
  \item 合格报告须含几何/材料、倒格截断与收敛、路径与模式追踪、频率/场型/对称标签、机器可读数据，以及传给 RCWA/CWT 的候选频窗和基底信息。
\end{enumerate}

\section{第 13 章}
\begin{enumerate}
  \item PWEM 给无损无限周期本征谱，RCWA 把层状周期单胞接到上下传播端口，可求反射、透射、吸收和开放共振；它仍未包含有限边缘与电热器件物理。
  \item 传输矩阵连续相乘会把指数增长与衰减的倏逝波混在一起而病态；散射矩阵通过端口入/出波并用稳定星积级联，避免巨大指数相消。
  \item RCWA 阈值增益是线性光学量；转为阈值电流还需材料增益--载流子关系、空间注入、复合、内部损耗和温度场。
  \item 固定 $\varepsilon'$ 且调 $\varepsilon''$ 忽略 Kramers--Kronig 联动、增益饱和、空间烧孔、载流子与热自洽，只表示线性极点到达实轴。
  \item 最小检查应含 Fourier 阶数/因子化、纵向切片、端口/PML、功率守恒、频率采样和模式追踪；每项需写明监视量及通过容差。
\end{enumerate}

\section{第 14 章}
\begin{enumerate}
  \item 有限阵列引入边缘泄漏、横向包络和尺寸离散模，可能改写单胞模式排序与远场，故需至少三种阵列尺寸及边界收敛。
  \item 宽带/瞬态通常 FDTD 有利；复杂孔形和局部加密通常 FEM 有利；高 $Q$ 可用 FEM 复本征值或 FDTD 长时谐波反演交叉验证；瞬态建立由 FDTD 直接给出。
  \item 网格：至少三档并跟踪复频率/模式重叠；边界：扫 PML 距离、厚度、吸收；提取：扫时间窗、监视点并用多指数残差。需预先声明误差预算。
  \item 示例：冷腔目标模 $Q$ 高且远场窄，但电流拥挤与热点偏向竞争模，使其净模增益先超过目标模；缺失电学、热学和有源竞争。
\end{enumerate}

\section{第 15 章}
\begin{enumerate}
  \item “准确”必须相对于同一边值问题和观测量；CWT、RCWA、FDTD、FEM 的几何、边界与状态空间不同，未指定问题时不能按单一数字排序。
  \item 一条合格路径：CWT/PWEM 排除错误频段与对称性；RCWA 检查端口与开放损耗；有限 FDTD/FEM 检查边缘、模式身份与远场；器件模型检查净增益和热工作点。每步淘汰标准需可量化。
  \item 例如 RCWA 中目标模最低辐射损耗，有限 FDTD 中却被边缘模超过；冲突来自有限终止/横向包络，而非必然来自某个求解器错误。也可给出几何或端口不一致反例。
  \item 绝对频率可受离散和边界共同偏移；模式排序、对称标签与趋势更能检验两方法是否识别同一支物理模式，但最终仍需绝对量的误差范围。
\end{enumerate}

\section{第 16 章}
\begin{enumerate}
  \item SCH 同时限制载流子和纵向光场，决定量子阱重叠、损耗层重叠与注入路径，是平面内光子晶体模型的纵向背景。
  \item 导带台阶限制电子外溢，价带台阶影响空穴注入与堆积；分配比例、温度和掺杂共同决定泄漏，不能只看总带隙差。
  \item $g_{\mathrm{mod}}=\Gamma_{\mathrm{opt}}g_{\mathrm{mat}}$ 增大同时，若 $\Gamma_{\mathrm{loss}}\alpha_{\mathrm{loss}}$ 增得更快，总阈值仍上升；还应提及注入/热代价。
  \item 失效场景：深刻蚀跨多层、上下强不对称、多个纵模近简并或金属引入强矢量边界；需三维全矢量开放模型。
\end{enumerate}

\section{第 17 章}
\begin{enumerate}
  \item MQW 可提高总模增益容量并降低单阱载流子要求，但可能加剧各阱注入不均、应变累积与热负荷；SQW 相反。三项都需比较。
  \item (001) 压应变通常抬升 HH 相对 LH、增强面内偶极 TE-like 增益；它与光子晶体的偏振损耗/辐射选择共同决定净优势，任一方不能单独保证输出偏振。
  \item 阻挡层差会使载流子泄漏或在空间上避开目标模高场区，$g_{\mathrm{mat}}(\rvec)$ 的加权积分下降；竞争模若与剩余增益区重叠更好会获益。
  \item 阻挡层改变泄漏、载流子分布、复合与压降，继而改变材料增益、模增益和阈值电流；“减少漏电”只描述其中一条路径。
\end{enumerate}

\section{第 18 章}
\begin{enumerate}
  \item 若 $L_s$ 远大于有源口径，横向载流子可近似铺平；若同量级或更小，电流拥挤和孔径边缘效应不可忽略。还需说明该判据只是首轮量级检查。
  \item 欧姆边界含有限接触比电阻与载流子交换；肖特基边界含势垒和偏压依赖。把非理想接触设为理想欧姆会低估压降、Joule 热和达到阈值载流子所需电流。
  \item 电流孔径改变 $N(\rvec)$ 与 $g_{\mathrm{mat}}(\rvec)$；各模按自身场权重积分，故相同总电流可给不同模增益并翻转竞争排序。
  \item 扩展层的折射率、自由载流子吸收、厚度和 Joule 热也改变光学模式与热透镜；因此电阻优化必须与损耗和模场联合。
\end{enumerate}

\section{第 19 章}
\begin{enumerate}
  \item 相同平均温升可对应不同热斑位置；模式按空间场权重感受 $\Delta n$、$\Delta g$、$\Delta\alpha$，局部热斑可重排竞争模而平均值看不出。
  \item 温升通过 $\partial n/\partial T$ 改变频率/模场，同时通过迁移率、接触和复合改变 $J,N$，再回写局域增益；答案需闭合到模式净增益。
  \item GaAs 与 GaN 的迁移率、掺杂激活、接触电阻、热导和工艺路径不同；应具体比较电流扩展与光学吸收/热代价，不能只写“材料不同”。
  \item 必须空间分辨的例子：毫米级口径含局部金属开窗和边缘散热；集中参数热阻无法给热点位置，也不能计算不同模的空间加权热漂移。
\end{enumerate}

\section{第 20 章}
\begin{enumerate}
  \item 单模稳定要求 $\Delta G_{\star j}$ 在整个电流、温度与容差范围内保持正值；目标模阈值低只说明某一点最先起振，不排除阈上交叉饱和与排序翻转。
  \item 边缘终止改变各横向包络的泄漏和相位，能有选择地提高高阶模损耗；因此它直接进入模鉴别矩阵，而非仅处理版图边界。
  \item 开放题须明确扰动尺度：单胞孔形优先改辐射矩阵；吸收环/终止优先改边缘损耗；电极/孔径则通过注入和温升改净模增益。若多条路径并存需说明主导依据。
  \item Lumerical 路径可用 MODE/RCWA $\to$ FDTD $\to$ CHARGE/HEAT；COMSOL 可用周期 Wave Optics $\to$ 有限 Wave Optics $\to$ Semiconductor/Heat Transfer。关键接口均包括模场、损耗、$N(\rvec)$、$T(\rvec)$ 与材料回写。
\end{enumerate}

\section{第 20B 章}
\begin{enumerate}
  \item 随机孔误差破坏相干叠加并混合邻近模式，远场旁瓣、偏振和模排序可在平均频率仅小幅变化时先退化；答案需提及空间相关长度。
  \item RCWA 回答单胞端口/损耗对深度的局部变化，不能给有限边缘；有限全波给模式与远场，不能给电流热回写；电热模型给工作点变化，仍需光学重算闭环。
  \item 计划至少含：按 wafer/lot 与局部随机量分层；用灵敏度筛关键参数；按联合分布/相关性采样；以明确观测量定义通过，并报告 $N_s$ 与置信区间。
  \item 金属与电流孔径错位同时改变吸收重叠、接触压降、电流密度和热点，进而通过 $\Delta n$、$\Delta g$ 和 $\Delta\alpha$ 回写模式；纯几何位移模型漏掉后两条链。
\end{enumerate}

\section{第 21 章}
\begin{enumerate}
  \item 观测时间越长，低频技术噪声、热漂移和模式跳变越多地进入谱；必须说明仪器分辨率、采样窗与内禀 Lorentzian/技术展宽的分离。
  \item 亮度与功率除以面积和立体角相关；功率增加若伴随发光面积、发散角或 $M^2$ 增大，亮度可不升反降。
  \item 可行链：竞争模获得净增益 $\to$ 输出在偏振/横向模间分配 $\to$ 偏振纯度下降、远场叠加变宽 $\to$ 模分配与频率噪声抬高线宽。
  \item 开放题按“求解器直接量—后处理公式—额外实验参数”评分。例如远场角谱可由全波直接得到，$M^2$ 需传播矩计算，工作线宽还需噪声模型与实验谱。
\end{enumerate}

\section{第 21B 章}
\begin{enumerate}
  \item 阈值是线性小信号交叉点；阈上还出现增益饱和、空间烧孔、多模交叉饱和、热漂移与外腔反馈，均可破坏单模稳态。
  \item 竞争模改变各 $S_m$ 及共享 $N$，横向/偏振场权重随之重分配，远场和偏振改变；模式分配噪声与幅相耦合又抬高线宽。
  \item 弛豫振荡位于快载流子--光子尺度，热漂移由热 RC 与边界决定；有限采样窗可把慢漂移看成随机噪声。建模应分别用速率方程与热方程并比较谱区。
  \item $\alpha_H$ 是载流子引起的折射率导数与增益导数之比，故材料/量子阱的色散与微分增益直接控制振幅--相位耦合；需注明波长、载流子和工作点。
  \item Purcell 因子还依赖模体积、发射体位置/谱失谐与单模弱耦合条件；大面积多通道 PCSEL 的 $V$ 和局域态密度并不由高 $Q$ 单独确定，$\beta$ 还取决于所有竞争辐射态。
\end{enumerate}

\section{第 22 章}
\begin{enumerate}
  \item L1 可报告被动候选频率/对称性；L2 可报告开放损耗与光学阈值排序；L3 加入输运后可估阈值电流来源；L4 加入热/动态/容差后才讨论工作窗口。每级需列缺失量。
  \item A 输出 $n_{\mathrm{eff}}$ 与重叠；B 输出候选频窗、端口和损耗；C 输出有限模式、远场与尺寸收敛；D 输出 $J,N,T$；E 输出回写后的净模裕量及不确定度。接口需保留单位、坐标和模式身份。
  \item Lumerical 的大规模有限 FDTD 常是算力瓶颈，跨 MODE/RCWA/FDTD/CHARGE 的几何与数据转换易错；COMSOL 的三维多物理耦合与非线性收敛可能最慢，材料/边界弱形式易错。允许不同答案但需具体。
  \item 停机判据示例：候选模对截断不收敛；有限尺寸排序随阵列/PML翻转；回写后净模裕量为负或不确定度大于裕量。三者分别阻止把数值伪影、单胞结论或名义点结论继续升级。
\end{enumerate}

\section{第 23 章}
\begin{enumerate}
  \item 在 $\{\bm u_{\mathrm B},\bm u_{\mathrm X},\ldots\}$ 基中给相应对角元分别加 $-\ii\gamma_{r,B}$、$-\ii\gamma_{r,X}$；若共用通道还应有 $-\ii d_Bd_X^*/2$ 的非对角项。物理上表示两模都向同一端口泄漏并可能耗散耦合。
  \item 由文中本征值显式代入 $\kappa_2\to-\kappa_2$ 比较实部；排序交换说明控制 $90^\circ$ 交叉耦合的相关几何 Fourier 分量改变符号/相位。完整分数需写出具体两个本征值。
  \item 汇总表至少列输入、求解模型、可测/可算输出、收敛证据、允许结论与缺失物理；不得把 PWEM/单胞结果写成阈值电流或稳定单模结论。
  \item 简化回写的材料、接触、热边界与增益模型尚未由独立实验标定；绝对 $I_{\mathrm{th}}$ 和线宽还受非辐射复合、技术噪声、模式竞争与封装影响。
\end{enumerate}

\section{第 24 章}
\begin{enumerate}
  \item 开放题按三步评分：先核对结论层级与所需变量；再核对几何、材料、单位、边界和模式身份；最后做网格/截断/迭代/后处理收敛。每步需指定可观察症状。
  \item 均匀增益忽略电流拥挤、泄漏、热斑与各模空间重叠，可能系统性低估阈值并错误排序竞争模；它只能作为线性光学代理。
  \item 最小对照可用同一材料平板或单孔周期单胞，统一坐标、单位、端口/PBC 相位和材料色散；比较解析频率/功率守恒及网格收敛，再逐项加回复杂结构。
  \item 清单至少覆盖共享输入、边界等价、模式追踪、功率守恒、网格/PML/阶次、求解器版本与误差预算；需把 Lumerical 与 COMSOL 的软件项映射到同一物理检查，而非并列按钮名。
\end{enumerate}
